% --- Organisations-Einheiten, Abteilungen & Begriffe (alphabetisch sortiert) ---

%%% Akronyme:
\myacronym{adac}{ADAC}{Allgemeiner Deutscher Automobil-Club e.V.}
\myacronym{ves}{VES}{ADAC Versicherung AG}
\myacronym{iag}{IAG}{Identity Access Governance}
\myacronym{refx}{REFX}{Flexible Real Estate Management}
\myacronym{po}{PO}{Product Owner}
\myacronym{dave}{DAVE}{Data Analytics in der Versicherung}
\myacronym{mvp}{MVP}{Minimal Viable Product}

\myacronym{cd}{FS-CD}{SAP Inkasso/Exkasso (FS-CD)}
\myacronym{tk}{TK}{Testkonzept}
\myacronym{sap}{SAP}{Systems, Applications, Products}[Integrierte Unternehmenssoftware zum \gls{erp}.]

\myacronym{insis}{Insis}{Insurance Information System}
\myacronym{ppa}{PPA}{Projektportfolio-Ausschuss}
\myacronym{bnc}{BnC}{\gls{insis}-Modul Billing and Collections}
\myacronym{ba}{BA}{Business-Analyst}
\myacronym{mce}{MCE}{Mitgliederservice \& Customer Experience}
\myacronym{ev}{ADAC e.V.}{ADAC e.V.}
\myacronym{sm}{SM}{Scrum Master}
\myacronym{vbb}{VBB}{Bestands- und Vertriebssysteme}
\myacronym{vss}{VSS}{Versicherung Support und Schaden}
\myacronym{wst}{WS Test}{Workstream Test}
\myacronym{scrum}{SCRUM}{Agiles Rahmenwerk zur iterativen Produktentwicklung}
\myacronym{vvq}{VVQ}{Bereich Vertrieb \& Querschnitt}
\myacronym{tm}{TM}{Testmanager}
\myacronym{sst}{SST}{Schnittstellen}

%% Versicherungsprodukte
\myacronym{rrv}{RRV}{ADAC Reiserücktritts-Versicherung}
\myacronym{unf}{UNF}{ADAC Unfallversicherung}
\myacronym{psv}{PSV}{ADAC Privatschutz – Kranken- und Pflegeversicherung}
\myacronym{aks}{AKS}{ADAC Auslandskrankenversicherung}
\myacronym{htr}{HTR}{ADAC Haustier-Reiseversicherung}
\myacronym{frv}{FRV}{ADAC Fahrrad-Versicherung}
\myacronym{saas}{SaaS}{Software-as-a-Service}
\myacronym{gu}{GU}{Generalunternehmer}
\myacronym{fb}{FB}{Fachbrereiche}
\myacronym{b2b}{B2B}{Business-to-Business}

\myacronym{epa}{EPA}{Produktivumgebung \gls{s4}}
\myacronym{zp1}{ZP1}{Produktivumgebung \gls{r3}}
\myacronym{dl}{DL}{Dienstleister}
\myacronym{bl}{BL}{Backlog}[Liste aller Aufgaben, Anforderungen und Arbeitspakete, die im Rahmen eines Projekts oder Programms zu erledigen sind. Das Backlog dient als zentrale Planungs- und Steuerungsgrundlage für die Projektarbeit.]



%%% Glossaries

\myacronym{adam20}{ADAM}{Allgemeiner Deutscher Automobilclub Mitgliedschaft 2.0}[Programm zur Ablösung der hostbasierten Mitgliedschaftsverwaltung des \gls{adac} inklusive Nebenbuch auf \gls{cd}-Basis.]

\myacronym{ricefw}{RICEFW}{Reports, Interfaces, Conversions, Enhancements, Forms, Workflows}[\gls{ricefw} bezeichnet im SAP-Kontext die sechs zentralen Entwicklungsobjekte – Reports (individuelle Auswertungen und Datendarstellungen), Interfaces (Schnittstellen zu Drittsystemen), Conversions (Programme zur Datenmigration und -formatkonvertierung), Enhancements (Erweiterungen der Standardfunktionalität), Forms (Druckausgaben und Layouts) und Workflows (automatisierte Geschäftsprozesse).]

\myacronym{vita}{VITA}{Versicherung IT des ADAC}[Transformationsprogramm zur Ablösung der hostbasierten Vertrags- und Schadensverwaltung der \gls{ves}.]

\myacronym{inex}{In/Ex}{In-/Exkasso}[Projekt unter \gls{vita} zur Implementierung eines neuen \gls{sl} der \gls{ves}.]

\myacronym{hypc}{HC}{HyperCare}[Phase der intensiven Betreuung direkt nach Produktivsetzungen, in der Stabilität, Fehlerbehebung und Support mit erhöhter Priorität sichergestellt werden, bevor der Regelbetrieb übernimmt.]

\myacronym{gl}{GL}{General Ledger}[Das \gls{gl} (Hauptbuch) ist das zentrale Register aller finanziellen Bewegungen eines Unternehmens und bildet die Grundlage des Jahresabschlusses.]

\myacronym{sl}{SL}{Sub-Ledger}[Das \gls{sl} (Nebenbuch) ist ein detailliertes Unterkonto des Hauptbuchs, das umfangreiche Transaktionen wie Debitoren- und Kreditorenprozesse abbildet und Genauigkeit, Effizienz und Revisionssicherheit erhöht.]

\myacronym{aif}{AIF}{Application Interface Framework}[Framework zur Entwicklung, Überwachung und Fehlerbehandlung von Schnittstellen im SAP-System.]

\myacronym{btp}{BTP}{Business Technology Platform}[Die \gls{btp} ist eine integrierte Technologiebasis von SAP zur Unterstützung digitaler Transformation.]

\myacronym{bfa}{BFA}{Brownfield Approach}[Transformationsansatz, bei dem bestehende Systeme, Prozesse und Daten weitgehend erhalten und in die neue Systemlandschaft integriert werden.]

\myacronym{adn}{ADN}{Aero-Dienst Nürnberg GmbH}[Tochtergesellschaft des ADAC und führender Dienstleister für Geschäfts- und Ambulanzluftfahrt in Europa.]

\myacronym{gfa}{GFA}{Greenfield Approach}[Vorgehensweise, bei der eine Lösung vollständig neu aufgebaut wird, im Gegensatz zum \gls{bfa}.]

\myacronym{op}{OP}{Offener Posten}[Ein offener Posten entsteht in der Nebenbuchhaltung, sobald eine Rechnung eingebucht, aber noch nicht ausgeglichen wurde. Erst durch Zahlungseingang oder -ausgang wird der Posten geschlossen.]

\myacronym{eol}{EOL}{End of Life}[Bezeichnet das Lebensende eines Produkts oder einer Software, wenn Herstellerunterstützung entfällt und dadurch Sicherheits- und Betriebsrisiken steigen.]

\myacronym{epic}{Epic}{Epic}[Eine umfangreiche Anforderung in der agilen Entwicklung, die in mehrere \gls{us} zerlegt wird.]

\myacronym{vbh}{VBH}{Bereich Schadenhilfe und -organisation}[Bereich der Versicherung mit IT-Koordination für die Systeme der Tätigen Hilfeorganisation, insbesondere das Softwareökosystem \gls{hermes}.]

\myacronym{hermes}{Hermes}{Hilfe Erbringung mit einem System}[System, das die gesamte Tätige Hilfeorganisation des ADAC abbildet – von Pannen- und Unfallhilfe bis zu weltweiten Krankenrücktransporten.]

\myacronym{e2e}{E2E}{End-to-End}[\gls{e2e} bezeichnet die durchgängige Ausführung eines Prozesses vom Auslöser bis zum Endergebnis ohne Unterbrechungen.]

\myacronym{f2g}{F2G}{Fit-Together}[Programm zur Vereinheitlichung der Systemumgebungs- und Testdatenlandschaft zur Absicherung der \gls{e2e}-Prozesse des \gls{adac}.]

\myacronym{voy}{Voyager}{Voyager}[Projekt zur Neugestaltung bzw. Ablösung des bisherigen Kampagnenmanagement-Tools.]

\myacronym{alina2}{Alina 2.0}{Alina 2.0}[Projekt zur Modernisierung des Vertriebsportals (Webseite und Verkaufsstellen inklusive Kassenanbindung).]

\myacronym{ais}{AIS}{ADAC IT Service GmbH}[Inhouse-IT-Dienstleister des gesamten ADAC-Verbunds.]
\myacronym{avm}{AVM}{ADAC Autovermietung GmbH}[Mietwagen- und Fuhrparkdienstleister des ADAC.]

\myacronym{bk}{BK}{Buchungskreis}[Kleinste organisatorische Einheit in SAP mit eigener Buchhaltung.]

\myacronym{nevesko}{NeVesKo}{Projekt neues Versicherungskonto}[Im Rahmen der 2014 durchgeführten Säulentrennung des \gls{adac} wurde die regulatorische Anforderung aufgestellt, die Versicherungsbeiträge künftig über ein der \gls{ves} gehörendes Bankkonto abzuwickeln. Bisher erfolgte dies über ein dem \gls{ev} gehörendes Bankverbindungskonto.]

\myacronym{tum}{T\&M}{Time and Material}[Vertragsmodell, bei dem die Abrechnung auf Basis der tatsächlich geleisteten Arbeitszeit und des Materialverbrauchs erfolgt, im Gegensatz zum Festpreis-Modell.]

\myacronym{cj}{CJ}{Customer Journey}[Modell zur Darstellung der Kundenreise durch verschiedene Touchpoints.]

\myacronym{br}{BR}{Betriebsrat}[Interessenvertretung der Beschäftigten.]

\myacronym{com}{COM}{ADAC Compliance Service GmbH}[Zentrale Compliance-Einheit des ADAC.]

\myacronym{con}{CON}{Controlling}[Abteilung für Planung, Steuerung, Analyse und Reporting im ADAC.]

\myacronym{dod}{DoD}{Definition of Done}[Abnahmekriterien in der agilen Produktentwicklung.]

\myacronym{dora}{DORA}{Digital Operational Resilience Act}[EU-Verordnung zur Stärkung der digitalen Resilienz.]

\myacronym{dp}{DP}{ADAC Datenportal}[Zentrales Datenmanagement-System.]

\myacronym{drc}{DRC}{Document and Reporting Compliance}[SAP-Modul für Dokumentations- und Berichtspflichten.]

\myacronym{dsgvo}{EU-DSGVO}{Datenschutz-Grundverordnung}[EU-weite Datenschutzregelung.]

\myacronym{erp}{ERP}{Enterprise Resource Planning}[Integrierte Software zur Steuerung von Unternehmensressourcen.]

\myacronym{f2h}{F2H}{Flight to HANA}[Programm zur S/4HANA-Einführung bei \gls{adn}.]

\myacronym{ft}{FT}{File Transfer}[Übertragung von Dateien zwischen Systemen oder Anwendungen über definierte Schnittstellen und Protokolle.]

\myacronym{fnz}{FNZ}{Abteilung Finanzen}[Bereich für Finanzplanung, Analyse und Reporting.]

\myacronym{gob}{GoB}{Grundsätze ordnungsmäßiger Buchführung}[Regelwerk für korrekte Buchführung.]

\myacronym{gobd}{GoBD}{Grundsätze ordnungsmäßiger Buchführung und Datenerfassung}[Digitale Buchführungsregeln.]

\myacronym{hr}{HR}{Human Resources}[Personalwesen.]

\myacronym{iea}{IEA}{Entwicklungsabteilung Enterprise Systeme}[Bereich für SAP-Entwicklung.]

\myacronym{ieb}{IEB}{Entwicklungsabteilung Enterprise Systeme}[Bereich für SAP-Betrieb und Architektur.]

\myacronym{is}{IS}{Integration Suite}[SAP-Integrationsplattform.]

\myacronym{itsec}{ITSEC}{IT-Sicherheit}[Stelle für IT-Sicherheitsmaßnahmen.]

\myacronym{itt}{ITT}{IT-Transformation Office}[Zentrale Abteilung für IT-Programmmanagement.]

\myacronym{jdr}{JDR}{Abteilung Digitalrecht}[Zentrale Stelle für digitalrechtliche Fragestellungen.]

\myacronym{jf}{JF}{Jour Fixe}[Regelmäßiger Besprechungstermin.]

\myacronym{jur}{JUR}{Juristische Zentrale}[Erstanlaufstelle für juristische Fragen.]

\myacronym{kpi}{KPI}{Key Performance Indicator}[Messgröße zur Bewertung von Leistung.]

\myacronym{kritis}{KRITIS}{Kritische Infrastrukturen}[Systeme mit Bedeutung für das Gemeinwohl.]

\myacronym{lns}{LnS}{Lift \& Shift}[Umzug der R/3-ECC-Landschaft in die Cloud.]

\myacronym{mbc}{MBC}{Multi-Bank Connectivity}[SAP-Modul zur Bankintegration.]

\myacronym{myit}{MyIT}{MyIT-Portal}[Internes Portal für Tickets und Services.]

\myacronym{pce}{PCe}{Private Cloud Edition}[Dedizierte SAP-Cloud-Lösung.]

\myacronym{pcr}{PCR}{Public Cloud Ready}[Architektur ohne spezifische Eingenporgrammierung nahe am SAP-Standard und Vorraussetzung für den Betrieb in einer SAP Public Cloud.]

\myacronym{pl}{PL}{Programmleiter}[Gesamtverantwortliche Person für ein Programm.]
\myacronym{tpl}{TPL}{Technische Programmleiterin}[Gesamtverantwortliche Person für die technischen Belange eines Programm.]
\myacronym{ag}{AG}{Auftraggeber / Projektsponsor}[Ein Projektauftraggeber / Projektsponsor ist eine Person oder eine Organisation, die ein Projekt initiiert und dessen Durchführung finanziert. Im \gls{adac} wird dies synonym verwendet.]

\myacronym{prjl}{PrjL}{Projektleitung}

\myacronym{pm}{PM}{Projektmanager}[Leitet und steuert Projekte.]

\myacronym{pmo}{PMO}{Program Management Office}[Zentrale Stelle für Programmmanagement.]

\myacronym{pmr}{PMR}{Projektmanagementreifegrad}[Grad der organisatorischen Projektmanagementfähigkeit.]

\myacronym{r2h}{R2H}{Road to HANA}[Transformationsprojekt zur S/4HANA-Einführung beim ADAC.]

\myacronym{r3}{R/3}{SAP ERP R/3}[Ältere Client-Server-ERP-Generation.]

\myacronym{raci}{RACI}{Responsibility Assignment Matrix}[Modell zur Rollen- und Verantwortlichkeitszuordnung.]

\myacronym{rasci}{RASCI}{Ressourcen, Verantwortlichkeiten, Zuständigkeiten, Informierte}[Erweitertes Rollenmodell.]

\myacronym{rec}{REC}{Rechtsabteilung}[Zentrale rechtliche Anlaufstelle.]

\myacronym{rv}{RV}{Rahmenvertrag}[Vertragliche Vereinbarung, die grundlegende Konditionen und Regelungen für wiederkehrende Leistungen oder Lieferungen festlegt.]

\myacronym{s4}{S/4}{SAP S/4HANA}[In-Memory-ERP-Generation.]

\myacronym{se}{SE}{ADAC Societas Europaea}[Europäische Aktiengesellschaft des ADAC.]

\myacronym{smartkpi}{SMART KPI}{SMART Key Performance Indicator}[Spezifisch, messbar, erreichbar, relevant und zeitgebunden.]

\myacronym{td}{TD}{Technical Debt}[Technische Schulden durch kurzfristige Lösungen.]

\myacronym{uac}{UAC}{Stonebranch Universal Automation Center}[Plattform für Automation und Orchestrierung.]

\myacronym{uat}{UAT}{User Acceptance Testing}[Abnahmetest mit Endanwendern.]
\myacronym{us}{US}{User Story}[Kurz formulierte, nutzerzentrierte Anforderung als Grundlage für Planung und Umsetzung in der agilen Produktentwicklung.]
\myacronym{sit}{SIT}{System Integrations Testing}

\myacronym{wp}{WP}{Wirtschaftsprüfer}[Unabhängige Prüfung der Finanzberichterstattung.]

\myacronym{ws}{WS}{Workstream}[Teilprojekt bzw. Arbeitsbereich.]

\myacronym{wsl}{WSL}{Workstream Lead}[Leitung eines Workstreams.]

\myacronym{stc}{STC}{Steering Committee}[Lenkungsausschuss für strategische Programmentscheidungen.]

\myacronym{wsr}{WSR}{WS Reports, Interfaces, Conversions, Enhancements, Forms, Workflows und Parallelbetrieb}[Workstream für RICEFW-Entwicklungsobjekte samt dem Parallelbetrieb.]

\myacronym{sod}{SOD}{Segregation of Duties}[Trennung von Verantwortlichkeiten zur Risikominimierung.]

\myacronym{zek}{ZEK}{Abteilung Zentraleinkauf}[Beschaffung von Waren und Dienstleistungen für den gesamten ADAC-Verbund.]

\myacronym{cop}{COP}{Cut Over Plan}[Ein Software Cut‑Over Plan ist ein strukturierter Ablaufplan, der alle notwendigen Schritte und Verantwortlichkeiten für den geordneten Übergang von einem alten zu einem neuen Softwaresystem während des Go‑Lives festlegt.]

\myacronym{cod}{COD}{Cut Over Durchführung}[Die Cut‑Over Durchführung bezeichnet die Phase, in der der geplante Übergang von einem alten zu einem neuen Softwaresystem tatsächlich umgesetzt wird, einschließlich aller technischen und organisatorischen Maßnahmen, um einen reibungslosen Wechsel sicherzustellen.]

\myacronym{pb}{PB}{Parallelbetrieb}[Der Parallelbetrieb ist eine Einführungsstrategie, bei der das alte und das neue System für einen Übergangszeitraum gleichzeitig genutzt werden, um Risiken zu minimieren und die korrekte Funktionsweise der neuen Software im produktiven Einsatz zu verifizieren.]

\myacronym{ere}{E-Invoice}{Elektronische Rechnungsstellung}[Die elektronische Rechnung ist eine Rechnung, die in einem strukturierten elektronischen Format ausgestellt, übermittelt und empfangen wird und automatisch maschinell verarbeitet werden kann – ohne Medienbrüche.]

\myacronym{cc}{CC}{Code Conversion}[Code Conversion ist der Vorgang, bestehenden Quellcode in ein anderes Format, eine andere Programmiersprache, eine neue Syntax oder eine modernisierte Struktur zu übertragen, um Kompatibilität, Wartbarkeit, Performance oder Systemmigration zu ermöglichen.]

\myacronym{cv}{CV}{Lebenslauf}

\myacronym{rup}{RUP}{Ramp-up-Phase}[Die Ramp-up-Phase bezeichnet im Projektmanagement den Zeitraum, in dem ein Projekt initial aufgebaut und die Ressourcen schrittweise in das Projekt integriert werden. In dieser Phase werden die Projektorganisation etabliert, die Prozesse definiert und die ersten Arbeitspakete umgesetzt, um das Projekt auf volle Fahrt zu bringen.]

\myacronym{nir}{NIR}{No Impact Rule}[Eine No Impact Rule (NIR) besagt, dass Änderungen oder Entwicklungen in einem Projekt keine negativen Auswirkungen auf bestehende Systeme, Prozesse oder Daten haben dürfen. Diese Regel dient dazu, die Stabilität und Integrität der bestehenden IT-Landschaft während der Implementierung neuer Lösungen sicherzustellen.]

\myacronym{it1}{IT1}{Integrationstest 1}[Ein Integrationstest, der die erste Phase der Integrationstests darstellt und darauf abzielt, die grundlegende Funktionalität und Interoperabilität der neu entwickelten oder angepassten Komponenten in einer kontrollierten Testumgebung zu überprüfen, bevor sie in die nächste Testphase übergehen.]